\documentclass{article}
\usepackage{amsmath,amssymb}
\usepackage{xurl}
\usepackage[hidelinks]{hyperref}
% Shared commands for notes in docs/paper-gaps/.

\DeclareUnicodeCharacter{2080}{\ifmmode{}_{0}\else\textsubscript{0}\fi}
\DeclareUnicodeCharacter{2081}{\ifmmode{}_{1}\else\textsubscript{1}\fi}
\DeclareUnicodeCharacter{2082}{\ifmmode{}_{2}\else\textsubscript{2}\fi}
\DeclareUnicodeCharacter{2099}{\ifmmode{}_{n}\else\textsubscript{n}\fi}

\newcommand{\Lean}{\textsc{Lean}}
\ExplSyntaxOn
\NewDocumentCommand{\leanid}{m}{
  \ifmmode
    \textsf{\detokenize{#1}}
  \else
    \group_begin:
    \urlstyle{sf}
    \tl_set:Nn \l_tmpa_tl {#1}
    \tl_replace_all:Nnn \l_tmpa_tl {\_} {_}
    \exp_args:NV \path \l_tmpa_tl
    \group_end:
  \fi
}
\ExplSyntaxOff
\providecommand{\tr}{\operatorname{tr}}
\newcommand{\tnleanIssueBase}{https://github.com/LionSR/TNLean/issues/}
\newcommand{\tnleanPrBase}{https://github.com/LionSR/TNLean/pull/}
\newcommand{\tnleanDocBase}{https://sirui-lu.com/TNLean/docs/}
\newcommand{\ghissue}[1]{\href{\tnleanIssueBase#1}{GitHub issue \##1}}
\newcommand{\ghpr}[1]{\href{\tnleanPrBase#1}{PR \##1}}
\newcommand{\doclink}[1]{\href{\tnleanDocBase#1.html}{\path{#1}}}

% Dirac notation and the identity operator, as in blueprint/src/macros/common.tex.
% \providecommand keeps note-local definitions authoritative.
\usepackage{dsfont}
\providecommand{\ket}[1]{|#1\rangle}
\providecommand{\bra}[1]{\langle #1|}
\providecommand{\braket}[2]{\langle #1 | #2 \rangle}
\providecommand{\ketbra}[2]{|#1\rangle\!\langle #2|}
\providecommand{\Id}{\mathds{1}}
\providecommand{\one}{\mathds{1}}
\providecommand{\C}{\mathbb{C}}
\providecommand{\MN}[1]{M_{#1}(\C)}
\providecommand{\MD}{M_D(\C)}

% Verdict marker: \gapnote{<kind>}{<status>}. Kind is the policy
% classification (clarification, local-correction, scope-restriction,
% unfaithful, false-source, open-gap); status is open, wip (elimination
% underway), resolved, or historical. Machine-read by the site index;
% prints a verdict line.
\newcommand{\gapnote}[2]{\par\noindent\textbf{Verdict:} #1 (#2).\par}

\title{A Corrected Approximating State for a Tensor with Repeated Blocks}
\author{}
\date{2026-10-01}
\begin{document}
\maketitle
\gapnote{local-correction}{resolved}

\section{What the source states}
Let $A^i=\bigoplus_{j=1}^b\operatorname{diag}(\mu_{j,1},\dots,\mu_{j,m_j})
    \otimes A_j^i$ with normal $A_j$ \cite[eq.~(S2)]{Malz2024Preparation}, so that
$\ket{\phi_N}\propto\sum_j\beta_j\ket{v_j}$ with
$\beta_j=\sum_{k=1}^{m_j}\mu_{j,k}^N$ and $\ket{v_j}$ the periodic state of
$A_j$ \cite[eqs.~(S3), (S4)]{Malz2024Preparation}. Block $q$ sites and write the
polar decomposition of the blocked map as $B=VP$. The source asserts
$P=\bigoplus_j\operatorname{diag}(\mu_{j,1}^q,\dots,\mu_{j,m_j}^q)\otimes P_j$
\cite[eq.~(S5)]{Malz2024Preparation}, replaces each $P_j$ by its fixed point,
and approximates the state on $N=qM$ sites by
\begin{align}
    \ket{\widetilde\phi_N}\propto V^{\otimes M}\sum_{j=1}^b\beta_j\ket{\Omega_j},
    \qquad \ket{\Omega_j}=\bigotimes_{k=1}^M\ket{\omega_j}_{R_kL_{k+1}}
    \label{eq:mswc24_rc_source_state}
\end{align}
\cite[eq.~(S7)]{Malz2024Preparation}. When some $m_j\geq2$ the block form fails
and no product of nearest-neighbour pairs with the weights $\beta_j$ gives an
approximation \cite{gap:mswc24_multiplicity_fixed_point}. This note gives the
corrected block form and the corrected approximating state, and shows that
part~(ii) of the approximation-error lemma holds for it when the $q$-site states
of distinct blocks are orthogonal. When they overlap, it proves the bound at the rate
of \cite{gap:mswc24_block_form_mixed_overlap} with the factor $\min(1,b)^{-1/2}$,
$b=\sum_j\lvert\beta_j\rvert^2$, which exceeds one when $b<1$; that open scope
restriction is recorded in \cite{gap:mswc24_repeated_overlap_small_weight_sum}.

\section{The corrected block form}
Let $E_{j,k}$ be the isometry placing the copy $k$ of block $j$ in the bond
space, so that $A^i=\sum_{j,k}\mu_{j,k}E_{j,k}A_j^iE_{j,k}^\dagger$ with
$E_{j,k}^\dagger E_{j',k'}=\delta_{jj'}\delta_{kk'}\mathbf 1$, and put
$K_{j,k}=E_{j,k}\otimes E_{j,k}$. For $q\geq1$ every word of length $q$ is a
direct sum, so the blocked map of $A$ is
\begin{align}
    B=\sum_{j,k}\mu_{j,k}^qB_jK_{j,k}^\dagger=\sum_jc_jB_jL_j^\dagger,
    \qquad
    c_j=\Bigl(\sum_k\lvert\mu_{j,k}\rvert^{2q}\Bigr)^{1/2},\qquad
    L_j=\sum_k\frac{\overline{\mu_{j,k}^q}}{c_j}K_{j,k},
    \label{eq:mswc24_rc_blocked}
\end{align}
with $B_j$ the blocked map of $A_j$, whenever some $\mu_{j,k}\neq0$. The maps
$L_j$ are isometries with orthogonal ranges, $L_j^\dagger L_{j'}=\delta_{jj'}$,
because $\sum_k\lvert\mu_{j,k}^q\rvert^2/c_j^2=1$. If the $q$-site states of
distinct blocks are orthogonal, $B_j^\dagger B_{j'}=0$ for $j\neq j'$, the polar
decomposition of an orthogonal sum \cite{gap:mswc24_block_form_mixed_overlap}
gives
\begin{align}
    V=\sum_jV_jL_j^\dagger,\qquad P=\sum_jc_jL_jP_jL_j^\dagger,
    \label{eq:mswc24_rc_polar}
\end{align}
with $B_j=V_jP_j$. On the copies of block $j$ the positive part is
$c_j\,w_jw_j^\dagger\otimes P_j$ with the unit vector
$w_{j,k}=\overline{\mu_{j,k}^q}/c_j$: it has rank one in the copy index, where
the source prints $\operatorname{diag}(\mu_{j,1}^q,\dots,\mu_{j,m_j}^q)$. For
$A^0=\operatorname{diag}(1,1,0)$, $A^1=\operatorname{diag}(0,0,1)$ this is the
matrix $\tfrac1{\sqrt2}\bigl(\begin{smallmatrix}1&1\\1&1\end{smallmatrix}\bigr)$
found in \cite{gap:mswc24_multiplicity_fixed_point}.

\section{The corrected approximating state}
Read $X=c_jL_jYL_j^\dagger$ as a tensor with physical leg the output pair and
virtual legs the input pair. Its entries are
$X^{(k'a',k'b')}_{(ka),(kb)}=c_j\,w_{j,k'}\overline{w_{j,k}}\,Y^{a'b'}_{ab}$ and
vanish off the copy-diagonal virtual pairs, so along a ring of $M$ sites the
virtual copy index is constant and
\begin{align*}
    \ket{\phi_M(X)}=c_j^M\Bigl(\sum_k\overline{w_{j,k}}^{\,M}\Bigr)
    L_j^{\otimes M}\ket{\phi_M(Y)}
    =\beta_jL_j^{\otimes M}\ket{\phi_M(Y)},
\end{align*}
since $c_j^M\overline{w_{j,k}}^{\,M}=\mu_{j,k}^{qM}$. With $Y=P_j$ this gives the
exact identity
$\ket{\phi_N}\propto V^{\otimes M}\sum_j\beta_jL_j^{\otimes M}\ket{v_{\mathrm{pos},j}}$,
the corrected form of the display after eq.~(S5) of the source, and replacing
$P_j$ by its fixed point gives the corrected approximating state
\begin{align}
    \ket{\widetilde\phi_N}\propto V^{\otimes M}\sum_j\beta_jL_j^{\otimes M}\ket{\Omega_j},
    \label{eq:mswc24_rc_state}
\end{align}
where $\ket{\Omega_j}$ is the product of the pairs $\omega_j$ of the fixed point of
$A_j$ in the bond space of block $j$. The weights $\beta_j$ are those of the
source; what changes is the fixed-point state: the isometry $L_j$ acts on the
two legs $L_kR_k$ of each site, so for $m_j\geq2$ the state
$L_j^{\otimes M}\ket{\Omega_j}$ carries the copy index on each site and is not a
product of nearest-neighbour pairs of the full bond space. It is still obtained
from a product of pairs by one on-site isometry per site. The states
$L_j^{\otimes M}\ket{\Omega_j}$ are orthonormal, so the fixed-point state has norm
$(\sum_j\lvert\beta_j\rvert^2)^{1/2}$ as in the source. For $m_j=1$ and
$\mu_j>0$, $L_j=K_j$ and (\ref{eq:mswc24_rc_state}) is the source's state
(\ref{eq:mswc24_rc_source_state}).

Since $VL_j=V_j$ by (\ref{eq:mswc24_rc_polar}),
\begin{align}
    V^{\otimes M}\sum_j\beta_jL_j^{\otimes M}\ket{\Omega_j}
    =\sum_j\beta_jV_j^{\otimes M}\ket{\Omega_j},
    \label{eq:mswc24_rc_after_v}
\end{align}
the combination, with the source's weights, of the approximating states of the
normal case for the blocks. Equivalently, the source's
state (\ref{eq:mswc24_rc_source_state}) with each pair on one copy
$k_j$ becomes correct after replacing $V$ by $V\sum_jL_jK_{j,k_j}^\dagger$.

\section{The error bound}
Assume $B_j^\dagger B_{j'}=0$ for $j\neq j'$ and $\beta\neq0$. This does not
force $\phi_N\neq0$, since a normal block can have $\phi_N(A_j)=0$ at small $N$
(for instance the Pauli matrices divided by $\sqrt3$ at $N=1$). The target is $\sum_j\beta_j\ket{v_j}$ and, by
(\ref{eq:mswc24_rc_after_v}), the corrected state is $\sum_j\alpha_j\ket{s_j}$ with
$\alpha_j=\beta_j/(\sum_l\lvert\beta_l\rvert^2)^{1/2}$ and
$\ket{s_j}=V_j^{\otimes M}\ket{\Omega_j}$. All cross overlaps vanish, as for
multiplicity one \cite{gap:mswc24_block_form_mixed_overlap}, so with
$p_j=\lvert\beta_j\rvert^2$, $z_j=\braket{s_j}{v_j}$ and $c_j'=\braket{v_j}{v_j}$,
\begin{align*}
    \lvert\braket{\widetilde\phi_N}{\phi_N}\rvert
    =\frac{\bigl\lvert\sum_jp_jz_j\bigr\rvert}
    {\bigl(\sum_jp_j\bigr)^{1/2}\bigl(\sum_jp_jc_j'\bigr)^{1/2}},
\end{align*}
and $\varepsilon\leq2\sum_j(\lvert z_j-1\rvert+\lvert c_j'-1\rvert)$. Part~(i) of the
lemma for each block bounds this by $Cy\,e^{Cy}$ with
$y=Me^{-\gamma q/\xi_{\mathrm{diag}}}$, which is $(N/q)e^{-\gamma q/\xi_{\mathrm{diag}}}$
for $q\geq1$, with no condition $q=o(N)$. The phases
and the possible cancellations in $\beta_j$ enter only through the weights $p_j$.

\section{Numerical check}
A script evaluates both states directly from the singular value decomposition of
$B$, for $N=qM$.
\begin{itemize}
    \item For $A^0=\operatorname{diag}(1,1,0)$, $A^1=\operatorname{diag}(0,0,1)$ the
          corrected state has overlap $1.000000$ with $\phi_N$ for $q=1,2$ and
          $M=1,2,3,4,6$, where the source's state with the pair on one copy gives
          $0.988496$, $0.948683$, $0.881546$, $0.800000$, $0.650791$, and the pair
          spread over both copies gives $1.000000$, $0.948683$, $0.800000$,
          $0.650791$, $0.502136$.
    \item For $d=4$, a block of bond dimension two on the physical indices
          $\{0,1\}$ with transfer eigenvalues of moduli $1$, $0.2421$, $0.2421$,
          $0.1929$ and copy weights $(1,0.8e^{0.7i})$, and a block of bond
          dimension one on $\{2,3\}$ with weight $0.9$, the corrected state gives
          $0.999321$, $0.999800$, $0.999644$ at $q=3$ and $M=1,2,3$, and
          $0.997257$, $0.995544$, $0.993003$, $0.990941$ at $q=2$ and $M=2,3,4,5$,
          where the source's state with the pair on one copy gives $0.997866$,
          $0.994090$, $0.992242$ and $0.983644$, $0.966207$, $0.954168$,
          $0.935826$. In every case the squared norm of the corrected fixed-point
          state equals $\sum_j\lvert\beta_j\rvert^2$.
    \item For the overlapping blocks $A_1=(1,0)$, $A_2=(3/5,4/5)$ of
          \cite{gap:mswc24_block_form_mixed_overlap} the corrected state is the
          source's state, and both give $0.857493$, $0.771528$, $0.713544$ at $q=1$
          and $M=2,4,6$: the correction concerns multiplicities, not overlaps.
\end{itemize}

\section{Blocks that occur}
The decomposition (S2) lists the normal tensors that occur in $A$, each copy
with its weight. A block with no copies, or a copy whose weight $\mu_{j,k}$
vanishes, is a summand with coefficient zero: it contributes nothing to $A$ and
only pads the bond dimension, and a block with no nonzero weight has $c_j=0$ and
$L_j$ undefined. The formalization adopts the intended reading as a convention of
the data: every block has at least one copy, and every copy has a nonzero
weight. This is a local fix of a degenerate reading, not a hypothesis added to
the source.

\section{What is formalized}
The corrected block form (\ref{eq:mswc24_rc_polar}), the identity
(\ref{eq:mswc24_rc_after_v}) and the overlap formula are formalized for arbitrary
multiplicities and nonzero complex weights, for blocks whose
$q$-site states are orthogonal, at every $q\geq1$; the overlap formula in addition
needs every blocked tensor $B_j$ to be injective, so that the vectors $\ket{s_j}$ are
unit vectors, and the error bound covers the shorter block lengths by
$\varepsilon\leq1$ after enlarging the constants. The bound
$\varepsilon\leq Cy\,e^{Cy}$, together with its $O$-form $\varepsilon\leq Cy$, is
formalized under the same assumptions, with $y=Me^{-\gamma q/\xi}$, at every $q\geq0$ and $M\geq1$ with
$\beta\neq0$.\footnote{Formal declarations:
    \leanid{MPSTensor.polarIso_blockTensor_repeatedBlockSum},
    \leanid{MPSTensor.polarPos_blockTensor_repeatedBlockSum} and
    \leanid{MPSTensor.copyApproxVector_repeatedBlockSum} in
    \doclink{TNLean/MPS/Preparation/RepeatedBlockSum};
    \leanid{MPSTensor.exists_approximationError_le_repeatedBlockSum} and
    \leanid{MPSTensor.exists_approximationError_le_mul_repeatedBlockSum} in
    \doclink{TNLean/MPS/Preparation/RepeatedBlockError}.} The orthogonality
restriction is the one recorded in \cite{gap:mswc24_block_form_mixed_overlap},
where for two overlapping blocks of multiplicity one the error lies between
$\tau^{2q}/16$ (for $N/q\geq3$) and $(N/q)\tau^{2q}$, with $\tau$ the eigenvalue of their mixed transfer
matrix, which the source's rate does not control; the construction here does not
change that example. For blocks of multiplicity one and weight one whose states
overlap, that note proves the bound with the correlation lengths of the mixed transfer
matrices included in the rate. The norm of the fixed-point state is formalized only as the
bound one after $V^{\otimes M}$; the overlap normalizes the state after $V^{\otimes M}$.

\section{Repeated blocks whose states overlap}
Drop the orthogonality of the $q$-site states of distinct blocks, let the blocks be normal in
the gauge $\sum_i(A_j^i)^\dagger A_j^i=\mathbf 1$ with positive definite fixed points
$\sigma_j$ of trace one, let $\lambda_2$ bound the moduli of the eigenvalues other than one of
every $E_{jj}$ and of all eigenvalues of the mixed transfer matrices $E_{jj'}$, $j\neq j'$, let
$\xi=-1/\ln\lvert\lambda_2\rvert$ and $0<\gamma<1/2$, and let $\lvert\mu_{j,k}\rvert\leq1$.
The corrected state (\ref{eq:mswc24_rc_state}) still converges, at the rate of
\cite{gap:mswc24_block_form_mixed_overlap}: there is $C$, independent of the weights, with
\begin{align}
    \varepsilon\leq\frac{Cy\,e^{Cy}}{\min(1,b)^{1/2}},\qquad y=Me^{-\gamma q/\xi},\qquad
    b=\sum_j\lvert\beta_j\rvert^2,
    \label{eq:mswc24_rc_overlap_bound}
\end{align}
for every $q\geq0$ and $M\geq1$ with $\beta\neq0$. For real weights $0\leq\mu_{j,k}\leq1$ of which
one equals one, $b\geq1$, and this is $\varepsilon\leq Cy\,e^{Cy}$ and $\varepsilon\leq C'y$.

The proof combines the two arguments. By (\ref{eq:mswc24_rc_blocked}) the blocked map is
$B=\sum_jc_jB_jL_j^\dagger$ with $c_j\leq m_j^{1/2}$, so the blocks of $B^\dagger B$ are
$c_jc_{j'}L_jB_j^\dagger B_{j'}L_{j'}^\dagger$, and the Gram estimate of
\cite{gap:mswc24_block_form_mixed_overlap}, which uses only that the embeddings have norm at most
one, gives $\lVert B^\dagger B-\sum_jc_j^2L_j(\sigma_j^T\otimes\mathbf 1)L_j^\dagger\rVert
\leq Ke^{-2\gamma q/\xi}$. The $\tfrac12$-H\"older bound of the square root gives
$\lVert P-P_\infty\rVert\leq K'e^{-\gamma q/\xi}$ with
$P_\infty=\sum_jc_jL_j(\sqrt{\sigma_j}^{\,T}\otimes\mathbf 1)L_j^\dagger$, whose square is the
limit of the Gram matrices since the $L_j$ are isometries with orthogonal ranges. Since
$V^\dagger B=P$, the overlap of the unnormalized state $V^{\otimes M}\sum_j\alpha_jL_j^{\otimes
M}\ket{\Omega_j}$ with $\ket{\phi_N(A)}$ is $b^{-1/2}\sum_j\overline{\beta_j}z_j$, with
$z_j=\braket{\Omega_j}{\phi_M(L_j^\dagger P)}$, where $L_j^\dagger P$ is read as a tensor whose
physical leg is a bond pair of block $j$. The row $L_j^\dagger P_\infty$ is
$\sum_k\mu_{j,k}^q(\sqrt{\sigma_j}^{\,T}\otimes\mathbf 1)K_{j,k}^\dagger$, and its mixed transfer
matrix against the fixed-point tensor of $\sigma_j$ is $\sum_k\mu_{j,k}^qR_k$ with
$R_k(\rho)=\operatorname{Tr}(E_{j,k}^\dagger\rho)\,E_{j,k}\sigma_j$. The $R_k$ are orthogonal
idempotents of trace one, so the powers $\sum_k\mu_{j,k}^{qn}R_k$ are bounded uniformly in the
weights and the $M$-th power has trace $\beta_j$; the telescoping estimate gives
$\lvert z_j-\beta_j\rvert\leq C_jy\,e^{C_jy}$. The states $L_j^{\otimes M}\ket{\Omega_j}$ are
orthonormal and $V^{\otimes M}$ does not increase norms, so normalizing the approximating state
only increases the overlap. Finally $\lvert\lVert\phi_N(A)\rVert^2-b\rvert\leq bK''e^{-\gamma
N/\xi}$, since every coefficient $\beta_j\overline{\beta_{j'}}$ has modulus at most $b$. With
$S=\sum_j\overline{\beta_j}z_j$ one has $\lvert S/b-1\rvert\leq b^{-1/2}\sum_j\lvert
z_j-\beta_j\rvert$, and the inequality of \cite{gap:mswc24_block_form_mixed_overlap} applied to
$S/b$ and $\lVert\phi_N(A)\rVert/b^{1/2}$ gives (\ref{eq:mswc24_rc_overlap_bound}).

The factor $\min(1,b)^{-1/2}$ is a scope restriction of the proof, not of the construction,
recorded in \cite{gap:mswc24_repeated_overlap_small_weight_sum}. It
exceeds one when $b<1$, for instance when the sums $\beta_j=\sum_k\mu_{j,k}^N$ cancel or every
weight has modulus below one, because the estimates of the $z_j$ are absolute; the source's bound
has no such factor. The numerical check below shows the error
converging at the same rate when $b$ is exponentially small. Removing the factor needs a
relative estimate of $z_j$ for blocks of small weight, the second-order perturbation bound for
the square root of $B^\dagger B$ that \cite{gap:mswc24_block_form_mixed_overlap} lists for the
sharp rate.

\section{Numerical check with overlaps}
A script evaluates the corrected state for two random normal blocks of bond dimension two on
$d=2$, with subleading eigenvalues $0.515$ and $0.850$ and mixed spectral radius $0.757$, so that
$\lvert\lambda_2\rvert=0.850$, from the eigendecomposition of $B^\dagger B$, using
$V^\dagger B=P$.
\begin{itemize}
    \item For copy weights $(1,0.8e^{0.7i})$ on the first block and $0.9$ on the second, and
          $M=8$, the error is $0.143$, $7.56\cdot10^{-3}$, $3.95\cdot10^{-4}$ at $q=4,8,12$, and
          $\varepsilon/(M\lvert\lambda_2\rvert^q)$ is $3.44\cdot10^{-2}$, $3.48\cdot10^{-3}$,
          $3.49\cdot10^{-4}$.
    \item For copy weights $(1,1)$ and $1$, and $M=8$, the corrected state gives $0.241$,
          $3.49\cdot10^{-2}$, $6.75\cdot10^{-3}$ at $q=4,8,12$, while the source's state with
          the pair on one copy gives $0.701$, $0.499$, $0.458$.
    \item For copy weights $(1,-1)$ and $0.5$ at odd $N$, where $\beta_1=0$, and $M=5$, the
          error is $0.424$, $7.54\cdot10^{-2}$, $1.51\cdot10^{-2}$ at $q=5,9,13$, while
          $b=0.5^{2N}$ is $8.9\cdot10^{-16}$, $8.1\cdot10^{-28}$, $7.3\cdot10^{-40}$.
\end{itemize}
The squared norm of the target is computed from
$\sum_{j,j'}\overline{\beta_j}\beta_{j'}\braket{\phi_N(A_j)}{\phi_N(A_{j'})}$; evaluating it as the
trace of the $N$-th power of the transfer matrix of $A$ loses all digits when $b$ is small.

\section{What is formalized with overlaps}
The bound (\ref{eq:mswc24_rc_overlap_bound}), and for real weights in $[0,1]$ one of which is one
its forms $\varepsilon\leq Cy\,e^{Cy}$ and $\varepsilon\leq C'y$, are formalized, with the
ingredients above.\footnote{Formal declarations:
    \leanid{MPSTensor.exists_approximationError_le_repeatedOverlappingBlockSum},
    \leanid{MPSTensor.exists_approximationError_le_repeatedOverlappingBlockSum_of_nonneg} and
    \leanid{MPSTensor.exists_approximationError_le_mul_repeatedOverlappingBlockSum_of_nonneg} in
    \doclink{TNLean/MPS/Preparation/RepeatedOverlappingBlockError};
    \leanid{MPSTensor.exists_norm_polarPos_blockTensor_repeatedBlockSum_sub_le},
    \leanid{MPSTensor.mixedMapLM_copyPosLimitRow_apply} and
    \leanid{MPSTensor.exists_norm_mpvOverlap_sub_bntWeight_le_of_norm_sub_copyPosLimitRow_le} in
    \doclink{TNLean/MPS/Preparation/RepeatedOverlappingBlockOverlap};
    \leanid{MPSTensor.exists_norm_gram_sum_sub_le} and
    \leanid{Matrix.norm_polarPos_sub_le_sqrt} in
    \doclink{TNLean/MPS/Preparation/OverlappingBlockGram}.}
The bound without the factor $\min(1,b)^{-1/2}$ when $b<1$ is not formalized
\cite{gap:mswc24_repeated_overlap_small_weight_sum}.

\section{Copies with weights of different phases}\label{sec:mswc24_rc_phases}
When a block has several copies whose weights have different phases, the phases cannot be
moved into the block as for multiplicity one \cite{gap:mswc24_block_form_mixed_overlap}, and
the question is which coefficients the approximating state needs. For the corrected state
(\ref{eq:mswc24_rc_state}) they are the coefficients $\beta_j=\sum_k\mu_{j,k}^N$ of the
source: the isometry $L_j$ carries the phases $\overline{\mu_{j,k}^q}/c_j$ of the copies, and
the bounds of the previous sections hold for complex weights. The sums
$\sum_k\lvert\mu_{j,k}\rvert^N$, which give the coefficient $\lvert\mu_j\rvert^N$ of
multiplicity one, and the copy norms $c_j^M$ do not work in their place. Take $d=2$, the
block $A_1=(1,0)$ with copy weights $(1,-1)$ and the block $A_2=(0,1)$ with weight one. For
odd $N$, $\beta=(0,1)$ and $\ket{\phi_N}\propto\ket{1\cdots1}$, and the corrected state with
the coefficients $(2,1)$ is proportional to $2\ket{0\cdots0}+\ket{1\cdots1}$, so
$\varepsilon=1-1/\sqrt5$ for every odd $q$ and $M$.

The source's construction, with the pairs of block $j$ on one copy $k_j$, also converges once
its coefficients are changed. For $q\geq1$ the blocked map is $B=\sum_jc_jB_jL_j^\dagger$
by (\ref{eq:mswc24_rc_blocked}), and $L_{j'}^\dagger K_{j,k}=\delta_{jj'}(\mu_{j,k}^q/c_j)$,
so $BK_{j,k}=\mu_{j,k}^qB_j=(\mu_{j,k}^q/c_j)BL_j$. The partial isometry $V$ vanishes on the
kernel of $B$: if $BY=0$ then $PY=0$, hence $\Pi Y=0$ for the support projector $\Pi$ of
$P$, and $VY=V\Pi Y=0$. Therefore
\begin{align}
    VK_{j,k}=\frac{\mu_{j,k}^q}{c_j}\,VL_j,
    \label{eq:mswc24_rc_one_copy}
\end{align}
for blocks whose $q$-site states overlap as well, and
$V^{\otimes M}K_{j,k_j}^{\otimes M}\ket{\Omega_j}=(\mu_{j,k_j}^q/c_j)^MV^{\otimes M}
    L_j^{\otimes M}\ket{\Omega_j}$. The source's state (\ref{eq:mswc24_rc_source_state}) with
the pairs on the copies $k_j$ and the coefficients
\begin{align}
    \beta'_j=\beta_j\Bigl(\frac{c_j}{\mu_{j,k_j}^q}\Bigr)^M
    \label{eq:mswc24_rc_one_copy_weight}
\end{align}
is the corrected state (\ref{eq:mswc24_rc_state}). For $m_j=1$, $c_j=\lvert\mu_j\rvert^q$ and
$\beta'_j=\lvert\mu_j\rvert^N$, the coefficient that repairs (S7) for multiplicity one
\cite{gap:mswc24_block_form_mixed_overlap}; for the tensor $A^0=\operatorname{diag}(1,1,0)$,
$A^1=\operatorname{diag}(0,0,1)$ of \cite{gap:mswc24_multiplicity_fixed_point},
$\beta'=(2\cdot2^{M/2},1)$, and the state is exact. The coefficients $\beta'_j$ depend on $q$
and $M$ separately, not only on $N$. With the printed coefficients $\beta_j$ the state
(\ref{eq:mswc24_rc_source_state}) fails, for $m_j\geq2$ already with positive weights
\cite{gap:mswc24_multiplicity_fixed_point}, and for $m_j=1$ as soon as a weight is not a
nonnegative real number \cite{gap:mswc24_block_form_mixed_overlap}.

\section{Numerical check with phases}
The script of the previous numerical checks evaluates both constructions for the two random
blocks used there, with copy weights $(1,e^{i})$ on the first block and $1$ on the second, at
$M=5$ and $q=3,5,7,9,11,13$.
\begin{itemize}
    \item The corrected state with $\beta_j=\sum_k\mu_{j,k}^N$ gives $0.477$, $9.75\cdot10^{-2}$,
          $9.74\cdot10^{-2}$, $1.67\cdot10^{-2}$, $9.35\cdot10^{-3}$, $6.94\cdot10^{-3}$.
    \item The corrected state with $\sum_k\lvert\mu_{j,k}\rvert^N$ gives $0.538$,
          $9.77\cdot10^{-2}$, $0.440$, $3.83\cdot10^{-2}$, $7.30\cdot10^{-2}$, $0.187$, and with
          $c_j^M$ it gives $0.549$, $0.114$, $0.558$, $7.73\cdot10^{-2}$, $0.122$, $0.261$.
    \item The source's state with the pairs on the first copy gives, with
          $\beta_j=\sum_k\mu_{j,k}^N$, $0.614$, $0.403$, $0.159$, $0.295$, $0.258$, $0.177$,
          and with $\sum_k\lvert\mu_{j,k}\rvert^N$, $0.609$, $0.403$, $0.195$, $0.291$,
          $0.254$, $0.188$.
    \item With the coefficients (\ref{eq:mswc24_rc_one_copy_weight}) the source's state gives
          the errors of the corrected state. Over three sets of weights, $q\in\{1,2,3,5\}$,
          $M\in\{1,3,4\}$ and both choices of the copy of the first block, the two errors agree
          to $6\cdot10^{-8}$, and (\ref{eq:mswc24_rc_one_copy}) holds to $10^{-14}$.
\end{itemize}

\section{What is formalized with phases}
The identity (\ref{eq:mswc24_rc_one_copy}), the equality of the overlaps of the source's state
with the coefficients (\ref{eq:mswc24_rc_one_copy_weight}) and of the corrected state, the value
$\beta'_j=\lvert\mu_j\rvert^N$ for a block with one copy, and the error bounds for the source's
state with these coefficients are formalized, for complex weights of any phase: for blocks
whose $q$-site states are orthogonal, $\varepsilon\leq Cy\,e^{Cy}$ and $\varepsilon\leq C'y$
at every $q\geq1$ and $M\geq1$ with $\beta\neq0$, and for overlapping blocks the bound
(\ref{eq:mswc24_rc_overlap_bound}) at every $q\geq1$, with its form for real weights in
$[0,1]$ one of which is one.\footnote{Formal declarations:
    \leanid{Matrix.polarIso_mul_eq_zero},
    \leanid{MPSTensor.polarIso_blockTensor_repeatedBlockSum_mul_pairEmbedding},
    \leanid{MPSTensor.oneCopyWeight},
    \leanid{MPSTensor.oneCopyWeight_of_subsingleton},
    \leanid{MPSTensor.nonNormalApproxOverlap_repeatedBlockSum_oneCopy},
    \leanid{MPSTensor.exists_approximationError_le_repeatedBlockSum_oneCopy},
    \leanid{MPSTensor.exists_approximationError_le_mul_repeatedBlockSum_oneCopy},
    \leanid{MPSTensor.exists_approximationError_le_repeatedOverlappingBlockSum_oneCopy} and
    \leanid{MPSTensor.exists_approximationError_le_repeatedOverlappingBlockSum_oneCopy_of_nonneg}
    in \doclink{TNLean/MPS/Preparation/OneCopyPairState}.} The failure of the coefficients
$\sum_k\lvert\mu_{j,k}\rvert^N$ and $c_j^M$ is checked numerically and not formalized.

\bibliographystyle{alpha}
\bibliography{references}
\end{document}
